\documentclass[12pt, letterpaper]{article}

% Packages for formatting
\usepackage[margin=1in]{geometry}
\usepackage{mathptmx} % Times New Roman font
\usepackage{setspace}
\singlespacing 
\usepackage{enumitem}
\usepackage{hyperref}
\usepackage{titlesec}
\usepackage{cite} % For handling citations

% Section formatting
\titleformat{\section}{\large\bfseries}{}{0em}{}
\titlespacing{\section}{0pt}{1.5ex plus 1ex minus .2ex}{1ex plus .2ex}

\begin{document}

\begin{center}
    \textbf{\Large Evaluating Small Fact-Checkers in Long-Context RAG via Adversarial Distractors} \\
    \vspace{0.5em}
    \textbf{Chenan Wang} \\
    \textbf{Project Type:} Research / Implementation
\end{center}

\vspace{1em}

\section*{B. Project Motivation}

\textbf{Problem Statement:} Current lightweight fact-checking models, such as MiniCheck \cite{tang2024minicheck}, are predominantly trained and evaluated on short, paragraph-level snippets (typically chunked to 350-500 tokens). However, in realistic Retrieval-Augmented Generation (RAG) pipelines, the context is often much longer and heavily heavily populated with irrelevant or conflicting information.

\textbf{Importance:} As LLM context windows grow, running fact-checking over massive retrieved documents becomes computationally expensive. While small fact-checkers (<1B parameters) offer a cost-effective alternative, their robustness against long contexts and adversarial distractors remains untested. Addressing this gap is critical for deploying reliable and affordable RAG systems.

\textbf{Context:} Recent studies indicate that LLMs frequently fail to utilize information buried in the middle of extended contexts \cite{laban2024summary}. It is currently an open question whether the local chunking mechanisms employed by models like MiniCheck \cite{tang2024minicheck} and AlignScore \cite{zha2023alignscore} can survive similar long-context stress tests without suffering significant accuracy degradation.

\textbf{Goal:} This project evaluates the MiniCheck model family (e.g., MiniCheck-FT5 and Bespoke-MiniCheck-7B) in long-context RAG scenarios. By injecting hard-negative adversarial distractors into the grounding documents, we aim to quantify the limits of small-parameter models in detecting hallucinations amidst noisy contexts.

\textbf{Prior Work:} As mentioned in context, current standardized factuality methods \cite{zha2023alignscore,tang2024minicheck} rely on fixed-length, relatively clean grounding documents. This project extends this paradigm by introducing dynamically noise-injected long contexts to expose the boundary failures of existing chunking strategies.

\section*{C. Milestones (Project Plan and Timeline)}

\begin{itemize}[leftmargin=*]
    \item \textbf{Week 1-2: Setup \& Baseline Replication} 
    \begin{itemize}
        \item \textit{Deliverable:} Configure local GPU environment; run the original MiniCheck codebase to replicate baseline results on the LLM-AGGREFACT benchmark.
    \end{itemize}
    \item \textbf{Week 3-4: Data Processing \& Noise Injection}
    \begin{itemize}
        \item \textit{Deliverable:} Build the long-context stress test suite. Sample subsets from LLM-AGGREFACT (e.g., WiCE, ExpertQA) and programmatically inject Wikipedia-sourced background noise and adversarial sentences into the grounding documents.
    \end{itemize}
    \item \textbf{Week 5-6: Experimentation}
    \begin{itemize}
        \item \textit{Deliverable:} Run inference across all test suites. Collect metrics on performance drops before and after the introduction of distractors.
    \end{itemize}
    \item \textbf{Week 7-8: Analysis \& Reporting}
    \begin{itemize}
        \item \textit{Deliverable:} Analyze why certain chunking mechanisms fail under long-context constraints; finalize the project report and presentation.
    \end{itemize}
\end{itemize}

\section*{D. Evaluation}

\textbf{Evaluation Approach:} An empirical ablation study. We will compare model performance on the original "clean" documents against the modified "noisy long documents" to measure robustness decay.

\textbf{Metrics/Criteria:} 
\begin{itemize}[leftmargin=*]
    \item \textit{Balanced Accuracy (BAcc):} To account for the natural class imbalance between factual and hallucinated samples.
    \item \textit{Inference Latency \& Cost:} To verify if the models remain computationally viable when processing extended contexts.
\end{itemize}

\textbf{Baselines and Comparisons:} 
The primary models will be MiniCheck-FT5 (770M) and Bespoke-MiniCheck-7B. AlignScore (355M) will serve as an earlier specialized baseline. A limited subset will also be evaluated using GPT-3.5/GPT-4 APIs to establish a performance upper bound.

\textbf{Data/Test Cases:} 
We will utilize the WiCE and ExpertQA subsets from the LLM-AGGREFACT benchmark, as they naturally require strong reasoning. Distractors will be sourced from random Wikipedia articles to significantly increase the token count and test the "needle in a haystack" capability.

\section*{E. Corresponding Resources}

\begin{itemize}[leftmargin=*]
    \item \textbf{Datasets/Corpora:} LLM-AGGREFACT benchmark; Wikipedia corpus (for distractors).
    \item \textbf{Tools/Libraries:} \texttt{Liyan06/MiniCheck} GitHub repository; Hugging Face \texttt{transformers} and \texttt{vLLM} (for 7B inference acceleration).
\end{itemize}

% Bibliography section
\bibliographystyle{unsrt}
\bibliography{refs}

\section*{F. Collaboration, Online Resources, and AI Usage Declaration}

\textbf{1. Collaboration:} 
I worked independently on this proposal and did not collaborate with other students.

\textbf{2. Online Resources:} 
\begin{itemize}[leftmargin=*]
    \item \url{https://github.com/Liyan06/MiniCheck}: Visited to understand the required environment and baseline capabilities of the MiniCheck models for the project feasibility check.
\end{itemize}

\textbf{3. AI Usage (gemini-pro-3):}
\begin{itemize}[leftmargin=*]
    \item \textbf{Goal:} I used gemini-pro-3 to help brainstorm the project direction regarding fact-checking in long-context RAG systems, translate my initial ideas into formal academic English, and format the proposal into standard LaTeX structure.
    \item \textbf{Summary of Learning:} Through the interaction, I refined the scope of my project in fact-checkers rather than a generic hallucination evaluation. I also learned how to structure a concise, one-page academic proposal using appropriate terminology.
    \item \textbf{Conversation Transcript:} The complete conversation transcript is attached as a separate file to this submission.
\end{itemize}

\end{document}